\chapter{Doctoral Entrance Problems}
\label{ch:problems}

The problems below are reproduced in their source language and routed by individual mathematical content.
\newpage

\section{École Nationale Supérieure de Mathématiques (ENSM)}
\begin{problem}{(ENSM) — Épreuve générale}{ra:ensm-2026-general-p1}
Let $\lambda\in]0,1[$ and let $f,g,h:[0,1]\to]0,+\infty[$ be continuous with
$$h(\lambda x+(1-\lambda)y)\ge f(x)^\lambda g(y)^{1-\lambda}.$$
Set $\alpha=\int_0^1f(x)\,dx$ and $\beta=\int_0^1g(x)\,dx$. Show that $\alpha,\beta\ne0$, that the normalized primitives are bijections, that $U=\lambda\Phi^{-1}+(1-\lambda)\Psi^{-1}$ is a $C^1$ bijection, and deduce
$$\int_0^1h(x)\,dx\ge\left(\int_0^1f(x)\,dx\right)^\lambda\left(\int_0^1g(x)\,dx\right)^{1-\lambda}.$$
\source{\emph{\small École Nationale Supérieure de Mathématiques (ENSM), 2026, Épreuve générale.}}\end{problem}
\newpage
\section{Université Djilali Liabès --- Sidi Bel Abbès}
\begin{problem}{(Sidi Bel Abbès) — Épreuve N°01}{ra:sba-2026-e1-p2}
On note $E(x)$ la partie entière. Montrer que pour tous $x,y\in\mathbb R$,
$$E(x)+E(y)+E(x+y)\le E(2x)+E(2y).$$
\source{\emph{\small Université Djilali Liabès - Sidi Bel Abbès, 2026, Épreuve~n\textsuperscript{o}~01.}}\end{problem}
\newpage
\section{Université des Sciences et de la Technologie Houari Boumediène (USTHB)}
\begin{problem}{(USTHB) — Épreuve N°02}{ra:usthb-2026-e2-p2}
Let $f:]0,1[\to\mathbb R$ be $f(x)=\sqrt{(1-x)/x}$. Prove it is bijective, find $f^{-1}$ and $(f^{-1})'(1)$, compute $\int 2y^2/(1+y^2)^2dy$, derive an integral formula for $f$, and solve
$$2x(1-x)y'+y=2(1-x)^2.$$
\source{\emph{\small Université des Sciences et de la Technologie Houari Boumediène (USTHB), 2026, Épreuve~n\textsuperscript{o}~02.}}\end{problem}
\newpage
\section{Université Kasdi Merbah --- Ouargla}
\begin{problem}{(Ouargla) — Épreuve N°01}{ra:ouargla-2025-e1-p1}
For $f(x)=\sum_{n\ge1}e^{-nx}/n^2$, study continuity on $[0,+\infty[$, $C^1$ and $C^2$ regularity on $\mathbb R_+^*$, differentiability at $0$, and prove
$$f(x)=f(0)+\int_0^x\ln(1-e^{-t})\,dt.$$
\source{\emph{\small Université Kasdi Merbah - Ouargla, 2025, Épreuve~n\textsuperscript{o}~01.}}\end{problem}
\newpage
\section{Université Abou Bekr Belkaïd --- Tlemcen}
\begin{problem}{(Tlemcen) — Épreuve N°02}{ra:tlemcen-2026-e2-p1}
For $\Phi(x)=\sum_{n\ge1}x^n/n^2$, determine the radius and boundary convergence, prove $\Phi(x)=\int_0^x-\ln(1-t)/t\,dt$, prove the reflection identity for $\Phi(x)+\Phi(1-x)$, and calculate $\Phi(1/2)$.
\source{\emph{\small Université Abou Bekr Belkaïd - Tlemcen, 2026, Épreuve~n\textsuperscript{o}~02.}}\end{problem}
\newpage
\section{Université Echahid Hamma Lakhdar --- El Oued}
\begin{problem}{(El Oued) — Épreuve N°01}{ra:eloued-2025-e1-p1}
Study continuity at $(0,0)$, partial and directional derivatives of
$$f(x,y)=\frac{x^2y+3y^3}{x^2+y^2},\quad f(0,0)=0,$$
and determine the Jacobian of $(f(x,y),f(y,x))$ at $(1,1)$.
\source{\emph{\small Université Echahid Hamma Lakhdar - El Oued, 2025, Épreuve~n\textsuperscript{o}~01.}}\end{problem}
\newpage
\section{Université Amar Telidji --- Laghouat}
\begin{problem}{(Laghouat) — Épreuve N°01}{ra:laghouat-2025-e1-p3}
State Beppo--Levi, prove the almost-everywhere convergence and integrability of a series $(f_n)$ under $\sum\int|f_n|<\infty$, and calculate
$$\int_0^1\frac{\ln x}{1-x}\,dx=-\frac{\pi^2}{6}.$$
\source{\emph{\small Université Amar Telidji - Laghouat, 2025, Épreuve~n\textsuperscript{o}~01.}}\end{problem}
\newpage
\section{Université Hassiba Benbouali --- Chlef}
\begin{problem}{(Chlef) — Épreuve N°01}{ra:chlef-2025-e1-p3}
Let $E=\mathbb R_1[x]$ and
$$f(p,q)=p(0)q(0)+\int_0^1p(t)q(1-t)\,dt.$$
Show that $Q(p)=f(p,p)$ is quadratic, that $f$ is an inner product, and find an orthonormal basis.
\source{\emph{\small Université Hassiba Benbouali - Chlef, 2025, Épreuve~n\textsuperscript{o}~01.}}\end{problem}
\newpage
\section{Université Djilali Liabès --- Sidi Bel Abbès}
\begin{problem}{(Sidi Bel Abbès) — Épreuve N°01}{ra:sba-2021-e1-p1}
Pour tout entier $n\in\mathbb N^*$, on pose
$$\mathbb J_n=\int_0^{+\infty}\frac{dx}{(1+x^3)^n}.$$
Montrer que $\mathbb J_n$ existe et est positive, établir une relation de récurrence entre $\mathbb J_n$ et $\mathbb J_{n+1}$, puis étudier la monotonie et la convergence de $(\mathbb J_n)$.
\source{\emph{\small Université Djilali Liabès - Sidi Bel Abbès, 2021, Épreuve~n\textsuperscript{o}~01.}}\end{problem}
